\section{The two-row ordinary label with a three-box tail}
\label{sec:alternating-three-row-tail}

\paragraph{Status and scope.}
New root derivation HD-AT301--03; the complete root author proof,
checker and first exact output have been read and passed.
The complete separate /root/g34\_eightfolds proof/source/code
audit and its distinct actual-subset reconstruction have also passed.
Shared preceding weighted-cycle authorship is explicitly disclosed.  No result in this section is promoted merely
because its finite tests pass.  The label is $(n-3,3)$.
Throughout $n\ge15$, $W=S^{(n-3,3)}$, and
\[
 a=\frac{n(n-1)(n-5)}6,\quad
 A=\frac{n^2-8n+3}6,\quad B=\frac{n(n-5)}6,\quad
 b=\frac n2,\quad d_0=A+B.                              \tag{AT3.1}
\]
The degree $a$ is an integer. The constants $A,B,b,d_0$ are
used as rational bounds and need not be integers.
For a permutation $x$, write $I(x)=n-\#\operatorname{cycles}(x)$,
including fixed letters.  Projective support is
$\nu(T)=\dim T-\max_z\dim\ker(T-zI)$.

\paragraph{The exact known character input.}
Let $P_j$ be the permutation representation on $j$-subsets.
The published two-row character formula gives
$\chi_W=\chi_{P_3}-\chi_{P_2}$, with the degree in (AT3.1).
We use James--Liebeck, \emph{Representations and Characters of
Groups}, second edition (2001), Theorem~29.13, pp.~344--345
\cite[Theorem~29.13, pp.~344--345]{HD:JamesLiebeck2001};
the entire proof and its preceding Proposition~29.12 were freshly
read in the institutional three-page scan.  This known result is
cited, not re-proved here.  Counting the fixed two- and three-subsets
therefore gives the elementary specialization
\[
 \chi_W(x)=\frac{f(f-1)(f-5)}6+(f-1)u+v,                 \tag{AT3.2}
\]
where $f,u,v$ are the numbers of one-, two-, and three-cycles of $x$.
The long label is not self-conjugate, so its restriction to $A_n$
is the ordinary irreducible module already used in the preceding
alternating-core interfaces.  Its rational $S_n$ extension supplies
all compatible outer automorphisms.

\subsection*{HD-AT301. A uniform index bound for every eigenvalue}

\begin{lemma}[HD-AT301]
For every nonidentity $x\in S_n$ and every scalar phase,
\[
                 \nu(Wx)\ge A I(x)+B\ge d_0,
 \qquad \chi_W(x)\ge-b.                                \tag{AT3.3}
\]
The character lower bound also holds at the identity.
\end{lemma}

\begin{proof}
If $f=0$, (AT3.2) is $-u+v\ge-n/2$.  If $f=1$, it is
nonnegative.  For $f=2,3,4$, its first term is respectively
$-1,-2,-2$, and its other terms are nonnegative.  For $f\ge5$
all terms are nonnegative.  This proves the character lower bound.

Put
\[
 c=\frac{n^2-6n+3}6,\qquad
 c_f=\frac{f^2+(n-9)f+(n-2)(n-4)}6.                     \tag{AT3.4}
\]
For $f\ge2$, the constraint $2u+3v\le n-f$ implies
$(f-1)u+v\le(f-1)(n-f)/2$.  Substitution in (AT3.2)
and exact division give, for $f<n$,
\[
             \frac{a-\chi_W(x)}{n-f}\ge c_f\ge c.
                                                                  \tag{AT3.5}
\]
The last inequality follows from
$6(c_f-c)=f^2+(n-9)f+5>0$.
For $f=0$, the character is at most $n/3$, giving the same
bound $c=a/n-1/3$.  For $f=1$, it is at most $(n-1)/3$,
and the corresponding quotient exceeds $c$ by $(n-5)/6$.
Thus $a-\chi_W(x)\ge c(n-f)$ for every $x$, including the
identity, where both sides are zero.

Let $L$ be the order of a fixed permutation $x$, let $m_z$
be the multiplicity of $z$ on $W$, and let $I=I(x)$.
Averaging over the cyclic group gives
\[
 a-m_1=\frac1L\sum_{j=0}^{L-1}(a-\chi_W(x^j))\ge cI.
                                                                  \tag{AT3.6}
\]
Here the average number of fixed letters of $x^j$ is the
number of cycles of $x$, by the actual permutation action.
Since the characters are real and at least $-b$, for $z\ne1$
with $z^L=1$ the exact Fourier formula yields
\[
 \begin{split}
 m_z-m_1
 &=\frac1L\sum_j\chi_W(x^j)(\operatorname{Re}z^{-j}-1)\\
 &\le\frac bL\sum_j(1-\operatorname{Re}z^{-j})=b.
 \end{split}                                                   \tag{AT3.7}
\]
The inequality $m_z\le m_1+b$ is also true for $z=1$.
Consequently $\nu(Wx)\ge cI-b$.

If $I\ge(n-2)/2$, then $c-A=n/3$ gives
\[
 \nu(Wx)-AI\ge\frac n3 I-\frac n2
                 \ge\frac{n(n-5)}6=B.                    \tag{AT3.8}
\]
Suppose instead $1\le I<(n-2)/2$.  A nontrivial cycle of
length $l$ moves $l\le2(l-1)$ letters, so every power of $x$
fixes at least $f_0=n-2I>2$ letters.  The polynomial $c_f$
is increasing for $f\ge0$, since $n\ge15$.  Applying (AT3.5)
to every nonidentity power improves (AT3.6) to
$a-m_1\ge c_{f_0}I$.  Together with (AT3.7), the desired
inequality is equivalent to
\[
 (n-f_0)\bigl(f_0^2+(n-9)f_0+2n+5\bigr)
                               \ge2n(n-2).                \tag{AT3.9}
\]
After subtracting the right side, its left side is
\[
 p(f_0)=-f_0^3+9f_0^2+(n^2-11n-5)f_0+9n.                 \tag{AT3.10}
\]
The polynomial is concave on $3\le f_0\le n-2$.
At $f_0=3$ and $n-2$ its values are respectively
$3n^2-24n+39$ and $2(n^2-11n+27)$, both positive for
$n\ge15$.  Hence it is positive throughout that interval.
For completeness the remaining integer endpoint $f_0=2$
has value $2n^2-13n+18>0$.
This proves (AT3.9), and therefore (AT3.3), for all nonidentity
permutations and all target eigenspaces.
\end{proof}

\subsection*{HD-AT302. The one-factor source, over every original field}

\begin{theorem}[HD-AT302]
In the actual rank-nine sixfold inducing-source setup, a single
ordinary core $A_n$ acting as $S^{(n-3,3)}$, $n\ge15$, is
impossible.  Every compatible scalar and outer automorphism is
allowed, and the original imaginary quadratic field is arbitrary.
\end{theorem}

\begin{proof}
The normalizer is scalar times the rational $S_n$ image.
That image has trivial scalar intersection: an element acting
as a scalar centralizes the faithful $A_n$ image, hence has
trivial conjugation on $A_n$; its $S_n$ permutation is the
identity.  The scalar and rational coordinates are therefore
actual homomorphisms.  The latter gives the genuine transitive
degree-$n$ quotient cover.

As in the preceding actual-source theorems, the Euler identity
$6=2a(h-1)+D$ excludes $h\ge2$, and $h=1$ has at most six
genuine marks and period rank at most six.  Rank nine forces
$h=0$, $r\ge12$, and $D=2a+6$.  The generating product-one
tuple in the actual $n$-letter quotient satisfies
$\sum I(x_i)\ge2n-2$ by Riemann--Hurwitz.
Every nonscalar genuine inertia has actual moving rank at
least $A I(x_i)+B$, by (AT3.3).  A genuine scalar has rank
$a\ge B$, while its index is zero, so the same formula applies.
It follows that
\[
 \begin{split}
 D-2a&\ge2A(n-1)+12B-2a\\
     &=n^2-8n-1>6 \qquad(n\ge15),                        \tag{AT3.11}
 \end{split}
\]
contradicting the Euler identity.  No conductor restriction
or equality of the original and inducing fields is used.
\end{proof}

\subsection*{HD-AT303. Every transitive tensor count over non-Eisenstein fields}

\begin{theorem}[HD-AT303]
Let the ordinary transitive core be $A_n^t$ on $W^{\otimes t}$,
where $W=S^{(n-3,3)}$, $n\ge15$ and $t\ge2$.
If the character field of the actual inducing representation
contains an imaginary quadratic field
$K\ne\mathbf Q(\sqrt{-3})$, it cannot give the actual
period-rank-nine sixfold source.  All compatible scalar,
outer and tensor-factor permutation actions are included.
\end{theorem}

\begin{proof}
Write $e=a^t$.  The same centralizer argument gives a unique
scalar times rational tensor-wreath normalizer, with trivial
scalar intersection.  Its rational coordinate gives the actual
transitive $nt$-letter action.  The scalar character has finite
image $\mu_M$, and the inducing character field is contained
in $\mathbf Q(\zeta_M)$.  Containment of the original
non-Eisenstein quadratic field forces at least two genuine
inertia marks with scalar outside $\mu_6$, by generation
and product one.  Equality of the two fields is not assumed.

The actual-cover argument again gives $h=0$, $r\ge12$,
$D=2e+6$.  The trace of every power of a rational tensor-wreath
operator is at least $-(b/a)e$: if the cycle-trace product
is negative, at least one factor has absolute value at most
$b$, and every other factor has absolute value at most $a$.
Rational cyclotomic multiplicities and the fourth-root square
trace estimate of D-ATT-02 consequently give moving rank
at least $(3-b/a)e/4$ at each of the two scalar marks.

Every genuine inertia moves at least $(d_0/a)e$ dimensions,
by (AT3.3) and convolution if it fixes factors, and by the
weighted-word bound if it moves them.  A factor-moving entry
has support at least $(a-1)e/(2a)$.  Here
$2d_0\le a-1$, $b\le d_0$, and $d_0\ge43$ for $n\ge15$;
these follow directly from (AT3.1).  If three entries move
factors, select the two scalar marks and a third distinct
factor-moving entry.  At least nine further genuine marks give
\[
 D\ge2e+a^{t-1}\left(9d_0-\frac{b+1}{2}\right)>2e+6.
                                                                  \tag{AT3.12}
\]
Thus exactly two entries move factors.  Product one makes their
factor permutations inverse, and transitivity makes each a full
$t$-cycle.  This is precisely the elementary two-moving-branch
argument of D-ATT-07 with the new constants.

Let $H$ be the sum of the natural indices of their two cycle
products.  Let $N$ be the total natural index of all within-factor
permutations at the other $r-2\ge10$ entries.  The actual
$nt$-letter quotient has total index
$2n(t-1)+H+N\ge2nt-2$, hence
\[
                              H+N\ge2n-2.                \tag{AT3.13}
\]
At a factor-fixing entry with $s\ge1$ active factors, convolution
and (AT3.3) give support at least the maximum of the $s$ active
supports.  It is therefore at least
$a^{t-1}(A I/s+B)\ge a^{t-1}(A I/t+B)$, where $I$ is that
entry's total natural index.  A genuine scalar has full rank
and satisfies the same bound with $I=0$, since $a\ge B$.

For a weighted full $t$-cycle with product $Wx$, the complete
general weighted-cycle lemma D-ATT-08 gives moving rank at least
\[
 \left(1-\frac1t\right)e
     +\frac{a^{t-1}}t\nu(Wx)-\frac{R_t(a)}t.              \tag{AT3.14}
\]
It remains valid when the product is the identity: here only
$\nu(Wx)\ge A I(x)$ is used, with both sides zero in that case.
The exact word error is $R_2(a)=a$; for $t\ge3$ the accepted
all-integer estimate is $R_t(a)\le2a^{t-2}$.  Thus
$2R_t(a)/a^{t-1}\le2$ for every $t\ge2$, because $a\ge350$.
Summing (AT3.14) and all factor-fixing estimates, using
(AT3.13), proves
\[
 \begin{split}
 D-2e
 &\ge\frac{a^{t-1}}t
       \left(2A(n-1)-2a+10tB-2\right)\\
 &=\frac{a^{t-1}}t\left(10tB-(n-1)^2-2\right)
   >B a^{t-1}>6.                                         \tag{AT3.15}
 \end{split}
\]
For the strict step, $9tB\ge18B=3n(n-5)$ and
$3n(n-5)-(n-1)^2-2=2n^2-13n-3>0$ for $n\ge15$.
This contradicts the actual Euler identity and proves the theorem.
\end{proof}

\paragraph{Strictness and the project frame.}
The earlier exact shape/conductor gate allows the label $(97,3)$
at $n=100$ over $\mathbf Q(i)$, for every positive tensor count:
its exact three-cycle fraction is below $1/10$ and its long row
exceeds the ALT07 threshold.  It is neither a standard nor a
two-box-tail label.  HD-AT302/303 remove that whole supplier
family, rather than restating the previous numerical cutoffs.
No actual source at that profile was previously asserted.
Other three-box and larger tails remain separate.

These are obstructions in the actual finite-cover representation
lambda-ring and its associated local systems; the scalar coordinate
and quotient cover are genuine homomorphisms and morphisms.
They shrink the condition-A supplier gate (cote16scan050;
SIGN_NEXT SN-C03/C04) and provide no condition-I cycle.
The preceding schemes-over-$\mathbf Z$, EGA/SGA/FGA, projector,
K-theoretic and motivic interpretation is unchanged.
No Markman, Severn or032 premise is added.
